% Notation
Dette projekt tager udgangspunkt i et fysisk system, der er beskrivet ved et koblet system af partielle differentialligninger.
For at holde rapporten notationsmæssigt konsistent bruges anvendes notationen inspireret af \cite{PDE_in_action}. 
For sæts og topology bruges:
\begin{gather*}
    \mathbb{N, Z, Q, R} \quad \text{(Naturlige, heltal, rationelle, reelle)}, \\
\end{gather*}

Generelt ville der bruges følgende notation for funktioner og rum:
\begin{table}[H]
    \centering
    \begin{tabular}{c|c|c}
    \text{Objekt} & \text{Notation} & \text{Eksempel} \\ \hline
    Skalar & Kursiv & $x, y, t, \alpha, \beta$ \\
    Vektor & Fed & $\mathbf{x, u, v}$ \\
    Matrice & Fed stor & $\mathbf{A, M, K}$ \\
    Funktioner & Kursiv & $f(\cdot), f_\theta, u, \sigma,$ \\
    Estimator & Hat& $\hat f_{theta}(\cdot), \hat{\mathbf{u}}_\theta$ \\
    \end{tabular}
\end{table}
Et objekt i n-dimensionelt rum skrives som $\mathbf{u}\in \mathbb{R}^n$. Disse bliver indekseret med $i, j, k$
For de givne operatorer kan der udføres forskellige operationer. Disse er noteret som følgene:

\begin{table}[H]
    \centering
    \begin{tabular}{c|c|c}
    \text{Operation} & \text{Notation} & \text{Eksempel} \\ \hline
    Transpose & $^T$ & $\mathbf{u}^T$ \\
    Inverse & $^{-1}$ & $\mathbf{A}^{-1}$ \\
    Middelværdi: $\mathbb{E}$ & Bar & $\overline{u}$ \\
    Varians: $\mathbb{V}$ & tilde & $\tilde{u}$ \\
    Afledte wrt. & $\partial$ & $\partial_x u$ \\
    Gradient & $\nabla$ & $\nabla u$ \\


    \end{tabular}
\end{table}



% 2.1.1 Funktioner og funktionrum
% 2.1.2 Differentialoperatorer
% 2.1.3 Ordinære differentialligninger
% 2.1.4 Partielle differentialligninger
% 2.1.5 Rand- og begyndelsesbetingelser
% 2.1.6 Approksimation af PDE-løsninger

Dette er også samlet i appendix A som en tabel

